\documentclass[12pt]{article}
\usepackage{amsmath,amssymb}

\begin{document}

\title{Reality Kernel Formalization}
\author{Gnome Badhi}
\date{22 May 2026}
\maketitle

\section*{Definitions}

Let
\[
\Sigma \text{ be the space of signals}, \quad
\mathcal{S} \text{ the space of states}, \quad
\mathcal{H} \text{ the space of integrity reports}.
\]

\subsection*{Sync frame}

The synchronization frame is
\[
\mathcal{F} = \bigl(f_{\text{Earth}},\, f_{\text{Magnetar}},\, H_0\bigr),
\]
where
\[
f_{\text{Earth}} = 7.83 \text{ Hz}, \quad
f_{\text{Magnetar}} \in \mathbb{R} \text{ (event date)}, \quad
H_0 = 73.2 \ \text{km/s/Mpc}.
\]

\subsection*{Operators}

\begin{itemize}
  \item Boundary operator:
  \[
  B : \Sigma \to \{0,1\}.
  \]

  \item Collapse operator:
  \[
  C : \Sigma \to \Sigma.
  \]

  \item Invariant set:
  \[
  \mathcal{I} = \{ I_1, I_2, \dots, I_n \}, \quad
  I_k : \mathcal{S} \times \Sigma \to \{0,1\}.
  \]

  \item Reporter:
  \[
  R : \mathcal{F} \times \{0,1\} \to \mathcal{H}.
  \]
\end{itemize}

\subsection*{Kernel}

The Reality Kernel is the tuple
\[
\mathcal{K} = \bigl(\mathcal{F}, B, C, \mathcal{I}, R\bigr).
\]

\section*{Processing Rule}

Given state \(s \in \mathcal{S}\) and signal \(\sigma \in \Sigma\):

\subsection*{Boundary check}

\[
B(\sigma) =
\begin{cases}
1, & \text{if the signal is within the accepted boundary},\\
0, & \text{otherwise}.
\end{cases}
\]

If \(B(\sigma) = 0\), the kernel returns a boundary rejection.

\subsection*{Invariant check}

All invariants must hold:
\[
\forall I_k \in \mathcal{I}:\ I_k(s,\sigma) = 1.
\]

If there exists \(k\) such that \(I_k(s,\sigma) = 0\), the kernel returns an invariant rejection.

\subsection*{Collapse}

If the boundary and all invariants are satisfied, the kernel applies the collapse operator:
\[
\text{out} = C(\sigma).
\]

Thus the full processing map is
\[
\text{Process}_{\mathcal{K}}(s,\sigma) =
\begin{cases}
\text{REJECT}_{\text{boundary}}, & B(\sigma) = 0, \\[4pt]
\text{REJECT}_{\text{invariant}}, & \exists k:\ I_k(s,\sigma) = 0, \\[4pt]
C(\sigma), & \text{otherwise}.
\end{cases}
\]

\section*{Integrity Rule}

Let \(\sigma_0 \in \Sigma\) be a neutral (no-op) probe signal.  
Define the integrity query:
\[
\text{Int}_{\mathcal{K}}(s) =
R\Bigl(\mathcal{F},\ \bigwedge_{k=1}^{n} I_k(s,\sigma_0)\Bigr),
\]
where
\[
\bigwedge_{k=1}^{n} I_k(s,\sigma_0) = 1
\quad \Leftrightarrow \quad
\forall k,\ I_k(s,\sigma_0) = 1.
\]

\end{document}
